% Options for packages loaded elsewhere
\PassOptionsToPackage{unicode}{hyperref}
\PassOptionsToPackage{hyphens}{url}
\documentclass[
]{article}
\usepackage{xcolor}
\usepackage[margin=1in]{geometry}
\usepackage{amsmath,amssymb}
\setcounter{secnumdepth}{-\maxdimen} % remove section numbering
\usepackage{iftex}
\ifPDFTeX
  \usepackage[T1]{fontenc}
  \usepackage[utf8]{inputenc}
  \usepackage{textcomp} % provide euro and other symbols
\else % if luatex or xetex
  \usepackage{unicode-math} % this also loads fontspec
  \defaultfontfeatures{Scale=MatchLowercase}
  \defaultfontfeatures[\rmfamily]{Ligatures=TeX,Scale=1}
\fi
\usepackage{lmodern}
\ifPDFTeX\else
  % xetex/luatex font selection
\fi
% Use upquote if available, for straight quotes in verbatim environments
\IfFileExists{upquote.sty}{\usepackage{upquote}}{}
\IfFileExists{microtype.sty}{% use microtype if available
  \usepackage[]{microtype}
  \UseMicrotypeSet[protrusion]{basicmath} % disable protrusion for tt fonts
}{}
\makeatletter
\@ifundefined{KOMAClassName}{% if non-KOMA class
  \IfFileExists{parskip.sty}{%
    \usepackage{parskip}
  }{% else
    \setlength{\parindent}{0pt}
    \setlength{\parskip}{6pt plus 2pt minus 1pt}}
}{% if KOMA class
  \KOMAoptions{parskip=half}}
\makeatother
\usepackage{color}
\usepackage{fancyvrb}
\newcommand{\VerbBar}{|}
\newcommand{\VERB}{\Verb[commandchars=\\\{\}]}
\DefineVerbatimEnvironment{Highlighting}{Verbatim}{commandchars=\\\{\}}
% Add ',fontsize=\small' for more characters per line
\usepackage{framed}
\definecolor{shadecolor}{RGB}{248,248,248}
\newenvironment{Shaded}{\begin{snugshade}}{\end{snugshade}}
\newcommand{\AlertTok}[1]{\textcolor[rgb]{0.94,0.16,0.16}{#1}}
\newcommand{\AnnotationTok}[1]{\textcolor[rgb]{0.56,0.35,0.01}{\textbf{\textit{#1}}}}
\newcommand{\AttributeTok}[1]{\textcolor[rgb]{0.13,0.29,0.53}{#1}}
\newcommand{\BaseNTok}[1]{\textcolor[rgb]{0.00,0.00,0.81}{#1}}
\newcommand{\BuiltInTok}[1]{#1}
\newcommand{\CharTok}[1]{\textcolor[rgb]{0.31,0.60,0.02}{#1}}
\newcommand{\CommentTok}[1]{\textcolor[rgb]{0.56,0.35,0.01}{\textit{#1}}}
\newcommand{\CommentVarTok}[1]{\textcolor[rgb]{0.56,0.35,0.01}{\textbf{\textit{#1}}}}
\newcommand{\ConstantTok}[1]{\textcolor[rgb]{0.56,0.35,0.01}{#1}}
\newcommand{\ControlFlowTok}[1]{\textcolor[rgb]{0.13,0.29,0.53}{\textbf{#1}}}
\newcommand{\DataTypeTok}[1]{\textcolor[rgb]{0.13,0.29,0.53}{#1}}
\newcommand{\DecValTok}[1]{\textcolor[rgb]{0.00,0.00,0.81}{#1}}
\newcommand{\DocumentationTok}[1]{\textcolor[rgb]{0.56,0.35,0.01}{\textbf{\textit{#1}}}}
\newcommand{\ErrorTok}[1]{\textcolor[rgb]{0.64,0.00,0.00}{\textbf{#1}}}
\newcommand{\ExtensionTok}[1]{#1}
\newcommand{\FloatTok}[1]{\textcolor[rgb]{0.00,0.00,0.81}{#1}}
\newcommand{\FunctionTok}[1]{\textcolor[rgb]{0.13,0.29,0.53}{\textbf{#1}}}
\newcommand{\ImportTok}[1]{#1}
\newcommand{\InformationTok}[1]{\textcolor[rgb]{0.56,0.35,0.01}{\textbf{\textit{#1}}}}
\newcommand{\KeywordTok}[1]{\textcolor[rgb]{0.13,0.29,0.53}{\textbf{#1}}}
\newcommand{\NormalTok}[1]{#1}
\newcommand{\OperatorTok}[1]{\textcolor[rgb]{0.81,0.36,0.00}{\textbf{#1}}}
\newcommand{\OtherTok}[1]{\textcolor[rgb]{0.56,0.35,0.01}{#1}}
\newcommand{\PreprocessorTok}[1]{\textcolor[rgb]{0.56,0.35,0.01}{\textit{#1}}}
\newcommand{\RegionMarkerTok}[1]{#1}
\newcommand{\SpecialCharTok}[1]{\textcolor[rgb]{0.81,0.36,0.00}{\textbf{#1}}}
\newcommand{\SpecialStringTok}[1]{\textcolor[rgb]{0.31,0.60,0.02}{#1}}
\newcommand{\StringTok}[1]{\textcolor[rgb]{0.31,0.60,0.02}{#1}}
\newcommand{\VariableTok}[1]{\textcolor[rgb]{0.00,0.00,0.00}{#1}}
\newcommand{\VerbatimStringTok}[1]{\textcolor[rgb]{0.31,0.60,0.02}{#1}}
\newcommand{\WarningTok}[1]{\textcolor[rgb]{0.56,0.35,0.01}{\textbf{\textit{#1}}}}
\usepackage{graphicx}
\makeatletter
\newsavebox\pandoc@box
\newcommand*\pandocbounded[1]{% scales image to fit in text height/width
  \sbox\pandoc@box{#1}%
  \Gscale@div\@tempa{\textheight}{\dimexpr\ht\pandoc@box+\dp\pandoc@box\relax}%
  \Gscale@div\@tempb{\linewidth}{\wd\pandoc@box}%
  \ifdim\@tempb\p@<\@tempa\p@\let\@tempa\@tempb\fi% select the smaller of both
  \ifdim\@tempa\p@<\p@\scalebox{\@tempa}{\usebox\pandoc@box}%
  \else\usebox{\pandoc@box}%
  \fi%
}
% Set default figure placement to htbp
\def\fps@figure{htbp}
\makeatother
\setlength{\emergencystretch}{3em} % prevent overfull lines
\providecommand{\tightlist}{%
  \setlength{\itemsep}{0pt}\setlength{\parskip}{0pt}}
\usepackage{bookmark}
\IfFileExists{xurl.sty}{\usepackage{xurl}}{} % add URL line breaks if available
\urlstyle{same}
\hypersetup{
  pdftitle={2501report\_ShenHongshan},
  pdfauthor={Shen Hongshan},
  hidelinks,
  pdfcreator={LaTeX via pandoc}}

\title{2501report\_ShenHongshan}
\author{Shen Hongshan}
\date{2025-11-20}

\begin{document}
\maketitle

\section{Predicting Mental Health Treatment Needs Based on Workplace
Factors and Demographic
Characteristics}\label{predicting-mental-health-treatment-needs-based-on-workplace-factors-and-demographic-characteristics}

\subsubsection{Author: Shen Hongshan
(3036290936)}\label{author-shen-hongshan-3036290936}

\subsubsection{\texorpdfstring{Data source:
\href{https://www.kaggle.com/datasets/osmi/mental-health-in-tech-survey/data}{Kaggle
Mental Health in Tech
Survey}}{Data source: Kaggle Mental Health in Tech Survey}}\label{data-source-kaggle-mental-health-in-tech-survey}

\subsection{\texorpdfstring{\textbf{1.
Introduction}}{1. Introduction}}\label{introduction}

\subsubsection{1.1 Background}\label{background}

Mental health issues in workplace environments represent a significant
public health concern with substantial economic and social consequences.
Early identification of individuals who need mental health support can
facilitate timely intervention, reduce workplace productivity loss, and
improve overall well-being. This project addresses the critical need for
data-driven approaches to mental health screening in professional
settings, particularly in high-pressure sectors like technology.

This project, analyzes the OSMI Mental Health in Tech Survey dataset
from Kaggle to model how workplace conditions and demographic factors
are associated with issue. The dataset contains survey responses to
workplace, demographic and mental health questions, providing rich
information for data-driven analysis of this spublic health issue.

\subsubsection{1.2 Data science question}\label{data-science-question}

The analysis focuses on two primary research questions:

\begin{itemize}
\item
  \textbf{What are the most significant demographics and workplace
  features associated with employees seeking mental health treatment?}
\item
  \textbf{Based on these associations, what targeted interventions can
  employers prioritize?}
\end{itemize}

\subsubsection{1.3 Methodological
Challenges}\label{methodological-challenges}

Challenges are anticipated in this analysis, including handling
\textbf{class imbalance} in the target variable, managing missing survey
data, and addressing potential \textbf{response bias}.The subjective
nature of self-reported mental health data may introduce measurement
noise, requiring careful data cleaning and validation approaches.

\subsection{\texorpdfstring{\textbf{2. Data Loading \&
Cleaning}}{2. Data Loading \& Cleaning}}\label{data-loading-cleaning}

\subsubsection{2.1 Data loading}\label{data-loading}

\begin{Shaded}
\begin{Highlighting}[]
\FunctionTok{library}\NormalTok{(tidyverse)}
\FunctionTok{library}\NormalTok{(readr)}
\FunctionTok{library}\NormalTok{(dplyr)}
\FunctionTok{library}\NormalTok{(ggplot2)}
\FunctionTok{library}\NormalTok{(knitr)}
\FunctionTok{library}\NormalTok{(caret) }
\FunctionTok{library}\NormalTok{(broom) }
\FunctionTok{library}\NormalTok{(car)  }\CommentTok{\# for VIF}
\FunctionTok{library}\NormalTok{(pROC) }
\FunctionTok{library}\NormalTok{(patchwork)}

\CommentTok{\# Read the dataset (make sure \textquotesingle{}survey.csv\textquotesingle{} is in the same folder)}
\NormalTok{survey\_data }\OtherTok{\textless{}{-}} \FunctionTok{read\_csv}\NormalTok{(}\StringTok{"survey.csv"}\NormalTok{)}

\CommentTok{\# Overview: dimensions and column names}
\FunctionTok{cat}\NormalTok{(}\StringTok{"Dataset dimensions:"}\NormalTok{, }\FunctionTok{dim}\NormalTok{(survey\_data), }\StringTok{"}\SpecialCharTok{\textbackslash{}n}\StringTok{"}\NormalTok{)}
\end{Highlighting}
\end{Shaded}

\begin{verbatim}
## Dataset dimensions: 1259 27
\end{verbatim}

\begin{Shaded}
\begin{Highlighting}[]
\FunctionTok{cat}\NormalTok{(}\StringTok{"Column names:}\SpecialCharTok{\textbackslash{}n}\StringTok{"}\NormalTok{)}
\end{Highlighting}
\end{Shaded}

\begin{verbatim}
## Column names:
\end{verbatim}

\begin{Shaded}
\begin{Highlighting}[]
\FunctionTok{names}\NormalTok{(survey\_data)}
\end{Highlighting}
\end{Shaded}

\begin{verbatim}
##  [1] "Timestamp"                 "Age"                      
##  [3] "Gender"                    "Country"                  
##  [5] "state"                     "self_employed"            
##  [7] "family_history"            "treatment"                
##  [9] "work_interfere"            "no_employees"             
## [11] "remote_work"               "tech_company"             
## [13] "benefits"                  "care_options"             
## [15] "wellness_program"          "seek_help"                
## [17] "anonymity"                 "leave"                    
## [19] "mental_health_consequence" "phys_health_consequence"  
## [21] "coworkers"                 "supervisor"               
## [23] "mental_health_interview"   "phys_health_interview"    
## [25] "mental_vs_physical"        "obs_consequence"          
## [27] "comments"
\end{verbatim}

\subsubsection{2.2 Handle Missing Values, drop Irrelevant biased sampled
variables}\label{handle-missing-values-drop-irrelevant-biased-sampled-variables}

\begin{itemize}
\item
  Check the number of missing value for each variable, delete those with
  too many NAs: comments, state
\item
  delete irrelevant variable for analysis: Timestamp
\item
  delete biased sampled variable: Country (60\% from USA, most countries
  which have only 1 respondents)
\end{itemize}

\begin{Shaded}
\begin{Highlighting}[]
\CommentTok{\# Missing value summary}
\FunctionTok{data.frame}\NormalTok{(}\AttributeTok{Missing\_data=}\FunctionTok{sort}\NormalTok{(}\FunctionTok{colSums}\NormalTok{(}\FunctionTok{is.na}\NormalTok{(survey\_data)),}\AttributeTok{decreasing=}\ConstantTok{TRUE}\NormalTok{))}
\end{Highlighting}
\end{Shaded}

\begin{verbatim}
##                           Missing_data
## comments                          1096
## state                              515
## work_interfere                     264
## self_employed                       18
## Timestamp                            0
## Age                                  0
## Gender                               0
## Country                              0
## family_history                       0
## treatment                            0
## no_employees                         0
## remote_work                          0
## tech_company                         0
## benefits                             0
## care_options                         0
## wellness_program                     0
## seek_help                            0
## anonymity                            0
## leave                                0
## mental_health_consequence            0
## phys_health_consequence              0
## coworkers                            0
## supervisor                           0
## mental_health_interview              0
## phys_health_interview                0
## mental_vs_physical                   0
## obs_consequence                      0
\end{verbatim}

\begin{Shaded}
\begin{Highlighting}[]
\CommentTok{\# Drop columns with too many NAs or irrelevant for analysis (comments, Timestamp, state)}
\NormalTok{survey\_clean0 }\OtherTok{\textless{}{-}}\NormalTok{ survey\_data }\SpecialCharTok{|\textgreater{}}
  \FunctionTok{select}\NormalTok{(}\SpecialCharTok{{-}}\NormalTok{Timestamp, }\SpecialCharTok{{-}}\NormalTok{Country, }\SpecialCharTok{{-}}\NormalTok{comments, }\SpecialCharTok{{-}}\NormalTok{state)}

\FunctionTok{cat}\NormalTok{(}\StringTok{"Columns left after removal:"}\NormalTok{, }\FunctionTok{names}\NormalTok{(survey\_clean0))}
\end{Highlighting}
\end{Shaded}

\begin{verbatim}
## Columns left after removal: Age Gender self_employed family_history treatment work_interfere no_employees remote_work tech_company benefits care_options wellness_program seek_help anonymity leave mental_health_consequence phys_health_consequence coworkers supervisor mental_health_interview phys_health_interview mental_vs_physical obs_consequence
\end{verbatim}

\subsubsection{2.3 Standardize Variables}\label{standardize-variables}

\begin{itemize}
\item
  Convert \textbf{Gender} variable into standard level: Male/ Female/
  Transgender/Non-binary/Others
\item
  Remove unrealistic \textbf{Age} input, only maintain age within
  18\textasciitilde100 years old
\item
  sort the levels for better visualization latter: work\_interfere(4
  level), leave(5 level), no\_employees(6 level)
\end{itemize}

\begin{Shaded}
\begin{Highlighting}[]
\NormalTok{survey\_clean0 }\OtherTok{\textless{}{-}}\NormalTok{ survey\_clean0 }\SpecialCharTok{|\textgreater{}}
  \FunctionTok{mutate}\NormalTok{(}
    \CommentTok{\# Gender standardization}
    \AttributeTok{Gender =} \FunctionTok{case\_when}\NormalTok{(}
      \FunctionTok{str\_detect}\NormalTok{(}\FunctionTok{tolower}\NormalTok{(Gender), }\StringTok{"\^{}m|\^{}male|\^{}man|\^{}maile|\^{}make|\^{}mal$|\^{}male{-}ish|\^{}guy"}\NormalTok{) }\SpecialCharTok{\textasciitilde{}} \StringTok{"Male"}\NormalTok{,}
      \FunctionTok{str\_detect}\NormalTok{(}\FunctionTok{tolower}\NormalTok{(Gender), }\StringTok{"\^{}f|\^{}female|\^{}woman|\^{}femaile|\^{}femake|\^{}fem"}\NormalTok{) }\SpecialCharTok{\textasciitilde{}} \StringTok{"Female"}\NormalTok{,}
      \ConstantTok{TRUE} \SpecialCharTok{\textasciitilde{}} \StringTok{"Transgender/Non{-}binary/Others"}
\NormalTok{    ),}
    \CommentTok{\# Convert character columns to factors}
    \FunctionTok{across}\NormalTok{(}\FunctionTok{where}\NormalTok{(is.character), as.factor)}
\NormalTok{  )}

\CommentTok{\# Handle age outliers}
\NormalTok{survey\_clean0 }\OtherTok{\textless{}{-}}\NormalTok{ survey\_clean0 }\SpecialCharTok{|\textgreater{}}
  \FunctionTok{mutate}\NormalTok{(}
    \AttributeTok{Age =} \FunctionTok{case\_when}\NormalTok{(}
\NormalTok{      Age }\SpecialCharTok{\textless{}} \DecValTok{18} \SpecialCharTok{\textasciitilde{}} \ConstantTok{NA}\NormalTok{,    }\CommentTok{\# Remove unrealistic ages}
\NormalTok{      Age }\SpecialCharTok{\textgreater{}} \DecValTok{100} \SpecialCharTok{\textasciitilde{}} \ConstantTok{NA}\NormalTok{,}
      \ConstantTok{TRUE} \SpecialCharTok{\textasciitilde{}}\NormalTok{ Age}
\NormalTok{    )}
\NormalTok{  ) }\SpecialCharTok{|\textgreater{}}
  \FunctionTok{filter}\NormalTok{(}\SpecialCharTok{!}\FunctionTok{is.na}\NormalTok{(Age))  }\CommentTok{\# Remove rows with invalid ages}

\CommentTok{\#sort the levels of variables}
\NormalTok{survey\_clean0}\SpecialCharTok{$}\NormalTok{work\_interfere }\OtherTok{\textless{}{-}} 
  \FunctionTok{factor}\NormalTok{(survey\_clean0}\SpecialCharTok{$}\NormalTok{work\_interfere,}\AttributeTok{levels =} \FunctionTok{c}\NormalTok{(}\StringTok{"Never"}\NormalTok{, }\StringTok{"Rarely"}\NormalTok{, }\StringTok{"Sometimes"}\NormalTok{, }\StringTok{"Often"}\NormalTok{))}
\NormalTok{survey\_clean0}\SpecialCharTok{$}\NormalTok{leave }\OtherTok{\textless{}{-}} 
  \FunctionTok{factor}\NormalTok{(survey\_clean0}\SpecialCharTok{$}\NormalTok{leave,}\AttributeTok{levels =} \FunctionTok{c}\NormalTok{(}\StringTok{"Very difficult"}\NormalTok{,}\StringTok{"Somewhat difficult"}\NormalTok{,}\StringTok{"Somewhat easy"}\NormalTok{,}\StringTok{"Very easy"}\NormalTok{,}\StringTok{"Don\textquotesingle{}t know"}\NormalTok{))}
\NormalTok{survey\_clean0}\SpecialCharTok{$}\NormalTok{no\_employees }\OtherTok{\textless{}{-}}
  \FunctionTok{factor}\NormalTok{(survey\_clean0}\SpecialCharTok{$}\NormalTok{no\_employees,}\AttributeTok{levels =} \FunctionTok{c}\NormalTok{(}\StringTok{"1{-}5"}\NormalTok{, }\StringTok{"6{-}25"}\NormalTok{, }\StringTok{"26{-}100"}\NormalTok{, }\StringTok{"100{-}500"}\NormalTok{, }\StringTok{"500{-}1000"}\NormalTok{, }\StringTok{"More than 1000"}\NormalTok{))}
\end{Highlighting}
\end{Shaded}

\subsubsection{2.3 Combine variables with similar
meaning}\label{combine-variables-with-similar-meaning}

\begin{itemize}
\tightlist
\item
  The variables ``benefits,'' ``care\_options,'' ``wellness\_program,''
  and ``seek\_help'' each measure a different facet of employers'
  support for mental health. So, they can be consolidated into a single
  variable, ``mental\_care,'' with a binary (Yes/No) level.
\end{itemize}

\begin{Shaded}
\begin{Highlighting}[]
\NormalTok{survey\_clean }\OtherTok{\textless{}{-}}\NormalTok{ survey\_clean0}\SpecialCharTok{|\textgreater{}}
  \FunctionTok{mutate}\NormalTok{(}\AttributeTok{mental\_care=}\FunctionTok{if\_any}\NormalTok{(}\FunctionTok{c}\NormalTok{(benefits, care\_options, wellness\_program, seek\_help), }\SpecialCharTok{\textasciitilde{}}\NormalTok{ . }\SpecialCharTok{==} \StringTok{"Yes"}\NormalTok{))}

\NormalTok{survey\_clean}\OtherTok{\textless{}{-}}\NormalTok{ survey\_clean}\SpecialCharTok{|\textgreater{}} 
  \FunctionTok{select}\NormalTok{(}\SpecialCharTok{{-}}\NormalTok{benefits, }\SpecialCharTok{{-}}\NormalTok{care\_options, }\SpecialCharTok{{-}}\NormalTok{wellness\_program, }\SpecialCharTok{{-}}\NormalTok{seek\_help)}
\end{Highlighting}
\end{Shaded}

\subsection{\texorpdfstring{\textbf{3. Exploratory Data Analysis \&
Visualization}}{3. Exploratory Data Analysis \& Visualization}}\label{exploratory-data-analysis-visualization}

\subsubsection{3.1 Target Variable Check}\label{target-variable-check}

\begin{itemize}
\item
  The target variable~\textbf{``treatment''}~indicates whether an
  employee is currently seeking treatment for a mental health condition.
\item
  \textbf{49.48\%}~report~\textbf{No}~and~\textbf{50.52\%}~report~\textbf{Yes},
  which underscoring the significance of mental health conditions in the
  workplace.
\end{itemize}

\begin{Shaded}
\begin{Highlighting}[]
\CommentTok{\#Distribution of treatment status}
\FunctionTok{as.data.frame}\NormalTok{(}\FunctionTok{table}\NormalTok{(survey\_clean}\SpecialCharTok{$}\NormalTok{treatment))}
\end{Highlighting}
\end{Shaded}

\begin{verbatim}
##   Var1 Freq
## 1   No  619
## 2  Yes  632
\end{verbatim}

\begin{Shaded}
\begin{Highlighting}[]
\CommentTok{\#Proportion table}
\FunctionTok{as.data.frame}\NormalTok{(}\FunctionTok{prop.table}\NormalTok{(}\FunctionTok{table}\NormalTok{(survey\_clean}\SpecialCharTok{$}\NormalTok{treatment)))}
\end{Highlighting}
\end{Shaded}

\begin{verbatim}
##   Var1      Freq
## 1   No 0.4948042
## 2  Yes 0.5051958
\end{verbatim}

\subsubsection{3.2 Demographic Overview}\label{demographic-overview}

\begin{itemize}
\tightlist
\item
  sample is skewed to male by Gender, thus it is important to notice the
  potential bias caused by this during analysis
\end{itemize}

\begin{Shaded}
\begin{Highlighting}[]
\NormalTok{demographic\_summary }\OtherTok{\textless{}{-}}\NormalTok{ survey\_clean }\SpecialCharTok{\%\textgreater{}\%}
  \FunctionTok{summarise}\NormalTok{(}
    \AttributeTok{n =} \FunctionTok{n}\NormalTok{(),}
    \AttributeTok{avg\_age =} \FunctionTok{mean}\NormalTok{(Age, }\AttributeTok{na.rm =} \ConstantTok{TRUE}\NormalTok{),}
    \AttributeTok{male\_pct =} \FunctionTok{round}\NormalTok{(}\FunctionTok{mean}\NormalTok{(Gender }\SpecialCharTok{==} \StringTok{"Male"}\NormalTok{)}\SpecialCharTok{*}\DecValTok{100}\NormalTok{,}\DecValTok{1}\NormalTok{),}
    \AttributeTok{female\_pct =} \FunctionTok{round}\NormalTok{(}\FunctionTok{mean}\NormalTok{(Gender }\SpecialCharTok{==} \StringTok{"Female"}\NormalTok{)}\SpecialCharTok{*}\DecValTok{100}\NormalTok{,}\DecValTok{1}\NormalTok{),}
    \AttributeTok{treatment\_rate =} \FunctionTok{round}\NormalTok{(}\FunctionTok{mean}\NormalTok{(treatment }\SpecialCharTok{==} \StringTok{"Yes"}\NormalTok{)}\SpecialCharTok{*}\DecValTok{100}\NormalTok{,}\DecValTok{1}\NormalTok{)}
\NormalTok{  )}
\CommentTok{\#Sample Demographic Summary}
\NormalTok{demographic\_summary}
\end{Highlighting}
\end{Shaded}

\begin{verbatim}
## # A tibble: 1 x 5
##       n avg_age male_pct female_pct treatment_rate
##   <int>   <dbl>    <dbl>      <dbl>          <dbl>
## 1  1251    32.1     78.7       19.8           50.5
\end{verbatim}

\subsubsection{3.3 Key Visualizations}\label{key-visualizations}

\begin{itemize}
\tightlist
\item
  Treatment rate by Age, Gender, family\_history, work\_interference,
  mental\_care, tech\_company
\end{itemize}

\begin{Shaded}
\begin{Highlighting}[]
\CommentTok{\# prepare Age{-}treatment information for visualization}
\NormalTok{Age\_summary}\OtherTok{\textless{}{-}}\NormalTok{survey\_clean}\SpecialCharTok{|\textgreater{}} 
  \FunctionTok{group\_by}\NormalTok{(Age)}\SpecialCharTok{|\textgreater{}}
  \FunctionTok{summarise}\NormalTok{(}
    \AttributeTok{n\_total =} \FunctionTok{n}\NormalTok{(),}
    \AttributeTok{n\_treatment =} \FunctionTok{sum}\NormalTok{(treatment }\SpecialCharTok{==} \StringTok{"Yes"}\NormalTok{, }\AttributeTok{na.rm =} \ConstantTok{TRUE}\NormalTok{),}
    \AttributeTok{prop\_treatment =}\NormalTok{ n\_treatment }\SpecialCharTok{/}\NormalTok{ n\_total,}
    \AttributeTok{.groups =} \StringTok{"drop"}
\NormalTok{  )}\SpecialCharTok{|\textgreater{}}
  \FunctionTok{filter}\NormalTok{(n\_total}\SpecialCharTok{\textgreater{}=}\DecValTok{5}\NormalTok{)}
\CommentTok{\#}
\NormalTok{Age\_visualization}\OtherTok{\textless{}{-}}\FunctionTok{ggplot}\NormalTok{(Age\_summary , }\FunctionTok{aes}\NormalTok{(}\AttributeTok{x =}\NormalTok{ Age, }\AttributeTok{y =}\NormalTok{ prop\_treatment, }\AttributeTok{size =}\NormalTok{ n\_total)) }\SpecialCharTok{+}
  \FunctionTok{geom\_point}\NormalTok{(}\AttributeTok{alpha =} \FloatTok{0.6}\NormalTok{) }\SpecialCharTok{+}
  \FunctionTok{geom\_smooth}\NormalTok{(}\AttributeTok{alpha =} \FloatTok{0.2}\NormalTok{)}\SpecialCharTok{+}
  \FunctionTok{theme\_minimal}\NormalTok{()}\SpecialCharTok{+}
  \FunctionTok{theme}\NormalTok{(}\AttributeTok{legend.position =} \StringTok{"none"}\NormalTok{)}

\NormalTok{Age\_visualization}
\end{Highlighting}
\end{Shaded}

\pandocbounded{\includegraphics[keepaspectratio]{2501report_ShenHongshan_Version2_files/figure-latex/unnamed-chunk-7-1.pdf}}

\begin{Shaded}
\begin{Highlighting}[]
\CommentTok{\# proportion of answers for different features }
\NormalTok{feature\_prop }\OtherTok{\textless{}{-}} \ControlFlowTok{function}\NormalTok{(data, var, var\_name) \{}
  \CommentTok{\# Use tidy evaluation with \{\{var\}\}}
\NormalTok{  percentage\_data }\OtherTok{\textless{}{-}}\NormalTok{ data }\SpecialCharTok{|\textgreater{}}
    \FunctionTok{count}\NormalTok{(\{\{var\}\}) }\SpecialCharTok{|\textgreater{}}
    \FunctionTok{mutate}\NormalTok{(}\AttributeTok{Percentage =}\NormalTok{ n }\SpecialCharTok{/} \FunctionTok{sum}\NormalTok{(n))}
  
  \CommentTok{\# Create plot}
  \FunctionTok{ggplot}\NormalTok{(percentage\_data, }\FunctionTok{aes}\NormalTok{(}\AttributeTok{x =}\NormalTok{ \{\{var\}\}, }\AttributeTok{y =}\NormalTok{ Percentage, }\AttributeTok{fill =}\NormalTok{ \{\{var\}\})) }\SpecialCharTok{+}
    \FunctionTok{geom\_bar}\NormalTok{(}\AttributeTok{stat =} \StringTok{"identity"}\NormalTok{) }\SpecialCharTok{+}
    \FunctionTok{scale\_y\_continuous}\NormalTok{(}\AttributeTok{labels =}\NormalTok{ scales}\SpecialCharTok{::}\FunctionTok{percent\_format}\NormalTok{(),}
                       \AttributeTok{limits =} \FunctionTok{c}\NormalTok{(}\DecValTok{0}\NormalTok{, }\FunctionTok{max}\NormalTok{(percentage\_data}\SpecialCharTok{$}\NormalTok{Percentage) }\SpecialCharTok{*} \FloatTok{1.15}\NormalTok{)) }\SpecialCharTok{+}
    \FunctionTok{labs}\NormalTok{(}\AttributeTok{title =} \FunctionTok{paste}\NormalTok{(var\_name, }\StringTok{"of Survey Respondents"}\NormalTok{),}
         \AttributeTok{x =}\NormalTok{ var\_name, }\AttributeTok{y =} \StringTok{"Percentage"}\NormalTok{) }\SpecialCharTok{+}
    \FunctionTok{theme\_minimal}\NormalTok{() }\SpecialCharTok{+}
    \FunctionTok{theme}\NormalTok{(}\AttributeTok{legend.position =} \StringTok{"none"}\NormalTok{)}
\NormalTok{\}}

\CommentTok{\# Treatment rate by feature (function) {-} FIXED to use tidy evaluation}
\NormalTok{treatment\_prop }\OtherTok{\textless{}{-}} \ControlFlowTok{function}\NormalTok{(data, feature, feature\_name) \{}
  \FunctionTok{ggplot}\NormalTok{(data, }\FunctionTok{aes}\NormalTok{(}\AttributeTok{x =}\NormalTok{ \{\{feature\}\}, }\AttributeTok{fill =}\NormalTok{ treatment)) }\SpecialCharTok{+}
    \FunctionTok{geom\_bar}\NormalTok{(}\AttributeTok{position =} \StringTok{"fill"}\NormalTok{) }\SpecialCharTok{+}
    \FunctionTok{labs}\NormalTok{(}\AttributeTok{y =} \StringTok{"Proportion"}\NormalTok{, }\AttributeTok{x =}\NormalTok{ feature\_name) }\SpecialCharTok{+}
    \FunctionTok{scale\_y\_continuous}\NormalTok{(}\AttributeTok{labels =}\NormalTok{ scales}\SpecialCharTok{::}\NormalTok{percent) }\SpecialCharTok{+}
    \FunctionTok{theme\_minimal}\NormalTok{()}
\NormalTok{\}}

\CommentTok{\#plots for each feature}
\FunctionTok{feature\_prop}\NormalTok{(survey\_clean, Gender, }\StringTok{"Gender"}\NormalTok{)}\SpecialCharTok{+}
  \FunctionTok{treatment\_prop}\NormalTok{(survey\_clean, Gender, }\StringTok{"Gender"}\NormalTok{)}
\end{Highlighting}
\end{Shaded}

\pandocbounded{\includegraphics[keepaspectratio]{2501report_ShenHongshan_Version2_files/figure-latex/unnamed-chunk-8-1.pdf}}

\begin{Shaded}
\begin{Highlighting}[]
\FunctionTok{feature\_prop}\NormalTok{(survey\_clean, family\_history, }\StringTok{"Family History"}\NormalTok{)}\SpecialCharTok{+}
  \FunctionTok{treatment\_prop}\NormalTok{(survey\_clean, family\_history, }\StringTok{"Family History"}\NormalTok{)}
\end{Highlighting}
\end{Shaded}

\pandocbounded{\includegraphics[keepaspectratio]{2501report_ShenHongshan_Version2_files/figure-latex/unnamed-chunk-8-2.pdf}}

\begin{Shaded}
\begin{Highlighting}[]
\FunctionTok{feature\_prop}\NormalTok{(survey\_clean, work\_interfere, }\StringTok{"Work Interference"}\NormalTok{)}\SpecialCharTok{+}
  \FunctionTok{treatment\_prop}\NormalTok{(survey\_clean, work\_interfere, }\StringTok{"Work Interference"}\NormalTok{)}
\end{Highlighting}
\end{Shaded}

\pandocbounded{\includegraphics[keepaspectratio]{2501report_ShenHongshan_Version2_files/figure-latex/unnamed-chunk-8-3.pdf}}

\begin{Shaded}
\begin{Highlighting}[]
\FunctionTok{feature\_prop}\NormalTok{(survey\_clean, tech\_company, }\StringTok{"Tech Company"}\NormalTok{)}\SpecialCharTok{+}
  \FunctionTok{treatment\_prop}\NormalTok{(survey\_clean, tech\_company, }\StringTok{"Tech Company"}\NormalTok{)}
\end{Highlighting}
\end{Shaded}

\pandocbounded{\includegraphics[keepaspectratio]{2501report_ShenHongshan_Version2_files/figure-latex/unnamed-chunk-8-4.pdf}}

\begin{Shaded}
\begin{Highlighting}[]
\FunctionTok{feature\_prop}\NormalTok{(survey\_clean, remote\_work, }\StringTok{"remote\_work"}\NormalTok{)}\SpecialCharTok{+}
  \FunctionTok{treatment\_prop}\NormalTok{(survey\_clean, remote\_work, }\StringTok{"remote\_work"}\NormalTok{)}
\end{Highlighting}
\end{Shaded}

\pandocbounded{\includegraphics[keepaspectratio]{2501report_ShenHongshan_Version2_files/figure-latex/unnamed-chunk-8-5.pdf}}

\begin{Shaded}
\begin{Highlighting}[]
\FunctionTok{feature\_prop}\NormalTok{(survey\_clean, self\_employed, }\StringTok{"self\_employed"}\NormalTok{)}\SpecialCharTok{+}
  \FunctionTok{treatment\_prop}\NormalTok{(survey\_clean, self\_employed, }\StringTok{"self\_employed"}\NormalTok{)}
\end{Highlighting}
\end{Shaded}

\pandocbounded{\includegraphics[keepaspectratio]{2501report_ShenHongshan_Version2_files/figure-latex/unnamed-chunk-8-6.pdf}}

\begin{Shaded}
\begin{Highlighting}[]
\CommentTok{\#company resourses}
\FunctionTok{feature\_prop}\NormalTok{(survey\_clean, mental\_care, }\StringTok{"Mental Health Care"}\NormalTok{)}\SpecialCharTok{+}
  \FunctionTok{treatment\_prop}\NormalTok{(survey\_clean, mental\_care, }\StringTok{"Mental Health Care"}\NormalTok{)}
\end{Highlighting}
\end{Shaded}

\pandocbounded{\includegraphics[keepaspectratio]{2501report_ShenHongshan_Version2_files/figure-latex/unnamed-chunk-8-7.pdf}}

\begin{Shaded}
\begin{Highlighting}[]
\FunctionTok{feature\_prop}\NormalTok{(survey\_clean, anonymity, }\StringTok{"anonymity for mental health information"}\NormalTok{)}\SpecialCharTok{+}
  \FunctionTok{treatment\_prop}\NormalTok{(survey\_clean, anonymity, }\StringTok{"anonymity for mental health information"}\NormalTok{)}
\end{Highlighting}
\end{Shaded}

\pandocbounded{\includegraphics[keepaspectratio]{2501report_ShenHongshan_Version2_files/figure-latex/unnamed-chunk-8-8.pdf}}

\begin{Shaded}
\begin{Highlighting}[]
\FunctionTok{feature\_prop}\NormalTok{(survey\_clean, leave, }\StringTok{"availability for taking medical leave for mental treatment"}\NormalTok{)}\SpecialCharTok{+}
  \FunctionTok{treatment\_prop}\NormalTok{(survey\_clean, leave, }\StringTok{"availability for taking medical leave for mental treatment"}\NormalTok{)}
\end{Highlighting}
\end{Shaded}

\pandocbounded{\includegraphics[keepaspectratio]{2501report_ShenHongshan_Version2_files/figure-latex/unnamed-chunk-8-9.pdf}}

\begin{Shaded}
\begin{Highlighting}[]
\FunctionTok{feature\_prop}\NormalTok{(survey\_clean, coworkers,}\StringTok{"feel free to discuss mental issue with coworkers"}\NormalTok{)}\SpecialCharTok{+}
  \FunctionTok{treatment\_prop}\NormalTok{(survey\_clean,coworkers,}\StringTok{"feel free to discuss mental issue with coworkers"}\NormalTok{)}
\end{Highlighting}
\end{Shaded}

\pandocbounded{\includegraphics[keepaspectratio]{2501report_ShenHongshan_Version2_files/figure-latex/unnamed-chunk-8-10.pdf}}

\subsubsection{\texorpdfstring{\textbf{3.4
Insights}}{3.4 Insights}}\label{insights}

\begin{itemize}
\tightlist
\item
  Age: Treatment needs show a continuous increase from age 18, peaking
  at around 40 years old, followed by a significant decline in later
  years. This pattern may reflect several factors: increased mental
  health awareness and openness to treatment in middle adulthood,
  workplace stressors reaching their peak during career-building years,
  or potentially changing attitudes toward mental health across
  generations.
\item
  Gender: Females show approximately 20\% higher treatment rates
  compared to males, while non-binary and gender-diverse individuals
  demonstrate over 80\% need for mental treatment.
\item
  family\_history: Those who have a family history of mental illness are
  around 35\% more likely to get treatment rather than without a family
  history. This aligns with clinical understanding, as family history is
  a well-established risk factor for many mental health disorders.
\item
  Work\_interfere: Respondents who report that mental health concerns
  ``often'' interfere with their work show the highest treatment rates,
  while those reporting ``never'' experiencing work interference show
  the lowest. This suggests that work interference serves as a
  significant motivator for seeking treatment, highlighting the
  \textbf{importance of workplace mental health interventions.}
\item
  tech\_company/ self\_employed: The little difference in treatment
  rates between tech company employees, non-tech employees, and
  self-employed individuals suggests that industry type and employment
  structure have limited influence on treatment-seeking behavior.
\item
  Mental healthcare/ anonymity/ leave: Companies that provide
  comprehensive mental health care programs---including confidential
  counseling services, robust privacy protections for mental health
  information, and clear mental health leave policies---show
  significantly higher treatment rates among employees. This supports
  the implementation of workplace mental health programs and policies
  are effective strategy for increasing treatment engagement.
\end{itemize}

\subsection{\texorpdfstring{\textbf{4. Modelling \& Statistical
Analysis}}{4. Modelling \& Statistical Analysis}}\label{modelling-statistical-analysis}

\subsubsection{4.1 Prepare Data for
Modelling}\label{prepare-data-for-modelling}

\begin{Shaded}
\begin{Highlighting}[]
\CommentTok{\# Create binary target variable}
\NormalTok{survey\_model }\OtherTok{\textless{}{-}}\NormalTok{ survey\_clean }\SpecialCharTok{\%\textgreater{}\%}
  \FunctionTok{mutate}\NormalTok{(}
    \AttributeTok{treatment\_binary =} \FunctionTok{ifelse}\NormalTok{(treatment }\SpecialCharTok{==} \StringTok{"Yes"}\NormalTok{, }\DecValTok{1}\NormalTok{, }\DecValTok{0}\NormalTok{)}
\NormalTok{  )}

\CommentTok{\# Select key predictors as per project proposal}
\NormalTok{predictors }\OtherTok{\textless{}{-}} \FunctionTok{c}\NormalTok{(}
  \StringTok{"Age"}\NormalTok{, }\StringTok{"Gender"}\NormalTok{, }\StringTok{"family\_history"}\NormalTok{,}
  \StringTok{"work\_interfere"}\NormalTok{, }\StringTok{"no\_employees"}\NormalTok{, }\StringTok{"remote\_work"}\NormalTok{, }\StringTok{"tech\_company"}\NormalTok{,}
  \StringTok{"mental\_care"}\NormalTok{,}
  \StringTok{"anonymity"}\NormalTok{, }\StringTok{"leave"}\NormalTok{,}
  \StringTok{"mental\_health\_consequence"}\NormalTok{, }\StringTok{"phys\_health\_consequence"}\NormalTok{,}
  \StringTok{"coworkers"}\NormalTok{, }\StringTok{"supervisor"}\NormalTok{, }\StringTok{"mental\_health\_interview"}\NormalTok{,}
  \StringTok{"mental\_vs\_physical"}\NormalTok{, }\StringTok{"obs\_consequence"}
\NormalTok{)}

\NormalTok{model\_data }\OtherTok{\textless{}{-}}\NormalTok{ survey\_model }\SpecialCharTok{|\textgreater{}}
  \FunctionTok{select}\NormalTok{(}\FunctionTok{all\_of}\NormalTok{(}\FunctionTok{c}\NormalTok{(}\StringTok{"treatment\_binary"}\NormalTok{, predictors))) }\SpecialCharTok{|\textgreater{}}
  \FunctionTok{na.omit}\NormalTok{() }\DocumentationTok{\#\#删除数据缺失行}

\FunctionTok{head}\NormalTok{(model\_data)}
\end{Highlighting}
\end{Shaded}

\begin{verbatim}
## # A tibble: 6 x 18
##   treatment_binary   Age Gender family_history work_interfere no_employees  
##              <dbl> <dbl> <fct>  <fct>          <fct>          <fct>         
## 1                1    37 Female No             Often          6-25          
## 2                0    44 Male   No             Rarely         More than 1000
## 3                0    32 Male   No             Rarely         6-25          
## 4                1    31 Male   Yes            Often          26-100        
## 5                0    31 Male   No             Never          100-500       
## 6                0    33 Male   Yes            Sometimes      6-25          
## # i 12 more variables: remote_work <fct>, tech_company <fct>,
## #   mental_care <lgl>, anonymity <fct>, leave <fct>,
## #   mental_health_consequence <fct>, phys_health_consequence <fct>,
## #   coworkers <fct>, supervisor <fct>, mental_health_interview <fct>,
## #   mental_vs_physical <fct>, obs_consequence <fct>
\end{verbatim}

\subsubsection{4.2 Split Data: Train/Test}\label{split-data-traintest}

\begin{Shaded}
\begin{Highlighting}[]
\FunctionTok{set.seed}\NormalTok{(}\DecValTok{42}\NormalTok{)}
\NormalTok{train\_idx }\OtherTok{\textless{}{-}} \FunctionTok{createDataPartition}\NormalTok{(model\_data}\SpecialCharTok{$}\NormalTok{treatment\_binary, }\AttributeTok{p=}\FloatTok{0.7}\NormalTok{, }\AttributeTok{list=}\ConstantTok{FALSE}\NormalTok{)}
\NormalTok{train\_data }\OtherTok{\textless{}{-}}\NormalTok{ model\_data[train\_idx, ]}
\NormalTok{test\_data }\OtherTok{\textless{}{-}}\NormalTok{ model\_data[}\SpecialCharTok{{-}}\NormalTok{train\_idx, ]}

\FunctionTok{cat}\NormalTok{(}\StringTok{"Training set size:"}\NormalTok{, }\FunctionTok{nrow}\NormalTok{(train\_data), }\StringTok{"}\SpecialCharTok{\textbackslash{}n}\StringTok{"}\NormalTok{)}
\end{Highlighting}
\end{Shaded}

\begin{verbatim}
## Training set size: 693
\end{verbatim}

\begin{Shaded}
\begin{Highlighting}[]
\FunctionTok{cat}\NormalTok{(}\StringTok{"Test set size:"}\NormalTok{, }\FunctionTok{nrow}\NormalTok{(test\_data), }\StringTok{"}\SpecialCharTok{\textbackslash{}n}\StringTok{"}\NormalTok{)}
\end{Highlighting}
\end{Shaded}

\begin{verbatim}
## Test set size: 296
\end{verbatim}

\subsubsection{4.3 Logistic Regression
Model}\label{logistic-regression-model}

target variable \& feature variable are both discrete, most of them with
only 2\textasciitilde3 levels. So, we can not use linear regression, but
logistic regression to explore what are the significant features
associated with our target value ``treatment''

\begin{Shaded}
\begin{Highlighting}[]
\NormalTok{logit\_model }\OtherTok{\textless{}{-}} \FunctionTok{glm}\NormalTok{(treatment\_binary }\SpecialCharTok{\textasciitilde{}}\NormalTok{ ., }\AttributeTok{data=}\NormalTok{train\_data, }\AttributeTok{family=}\StringTok{"binomial"}\NormalTok{)}

\FunctionTok{cat}\NormalTok{(}\StringTok{"Model summary (logistic regression):}\SpecialCharTok{\textbackslash{}n}\StringTok{"}\NormalTok{)}
\end{Highlighting}
\end{Shaded}

\begin{verbatim}
## Model summary (logistic regression):
\end{verbatim}

\begin{Shaded}
\begin{Highlighting}[]
\FunctionTok{summary}\NormalTok{(logit\_model)}
\end{Highlighting}
\end{Shaded}

\begin{verbatim}
## 
## Call:
## glm(formula = treatment_binary ~ ., family = "binomial", data = train_data)
## 
## Coefficients:
##                                     Estimate Std. Error z value Pr(>|z|)    
## (Intercept)                         -3.05835    0.89832  -3.405 0.000663 ***
## Age                                  0.03011    0.01450   2.077 0.037837 *  
## GenderMale                          -0.79747    0.28035  -2.845 0.004448 ** 
## GenderTransgender/Non-binary/Others  0.19497    0.91680   0.213 0.831584    
## family_historyYes                    1.08318    0.21276   5.091 3.56e-07 ***
## work_interfereRarely                 2.11481    0.32405   6.526 6.75e-11 ***
## work_interfereSometimes              2.68910    0.29640   9.072  < 2e-16 ***
## work_interfereOften                  2.88780    0.38183   7.563 3.94e-14 ***
## no_employees6-25                     0.36963    0.35319   1.047 0.295310    
## no_employees26-100                   0.65790    0.36357   1.810 0.070367 .  
## no_employees100-500                  0.70573    0.43058   1.639 0.101205    
## no_employees500-1000                 0.85216    0.60058   1.419 0.155930    
## no_employeesMore than 1000           0.26291    0.37641   0.698 0.484892    
## remote_workYes                       0.01822    0.23748   0.077 0.938838    
## tech_companyYes                     -0.06577    0.28103  -0.234 0.814958    
## mental_careTRUE                      0.57800    0.22927   2.521 0.011699 *  
## anonymityNo                         -0.01578    0.45670  -0.035 0.972441    
## anonymityYes                         0.47029    0.26063   1.804 0.071163 .  
## leaveSomewhat difficult             -0.27922    0.47677  -0.586 0.558115    
## leaveSomewhat easy                  -0.72958    0.44131  -1.653 0.098288 .  
## leaveVery easy                      -0.32236    0.48513  -0.664 0.506384    
## leaveDon't know                     -0.51625    0.41638  -1.240 0.215023    
## mental_health_consequenceNo         -0.25477    0.28531  -0.893 0.371878    
## mental_health_consequenceYes        -0.11745    0.30379  -0.387 0.699050    
## phys_health_consequenceNo            0.17988    0.27915   0.644 0.519329    
## phys_health_consequenceYes          -0.61063    0.48552  -1.258 0.208507    
## coworkersSome of them                0.44016    0.28520   1.543 0.122747    
## coworkersYes                         1.09031    0.42254   2.580 0.009870 ** 
## supervisorSome of them              -0.25823    0.29497  -0.875 0.381345    
## supervisorYes                       -0.32457    0.33779  -0.961 0.336621    
## mental_health_interviewNo            0.31513    0.31713   0.994 0.320365    
## mental_health_interviewYes           0.16143    0.72359   0.223 0.823457    
## mental_vs_physicalNo                -0.04933    0.26430  -0.187 0.851929    
## mental_vs_physicalYes                0.04360    0.29458   0.148 0.882340    
## obs_consequenceYes                   0.23140    0.30289   0.764 0.444886    
## ---
## Signif. codes:  0 '***' 0.001 '**' 0.01 '*' 0.05 '.' 0.1 ' ' 1
## 
## (Dispersion parameter for binomial family taken to be 1)
## 
##     Null deviance: 903.92  on 692  degrees of freedom
## Residual deviance: 622.27  on 658  degrees of freedom
## AIC: 692.27
## 
## Number of Fisher Scoring iterations: 5
\end{verbatim}

\paragraph{Predictors sorted by p-value
(significance)}\label{predictors-sorted-by-p-value-significance}

\begin{Shaded}
\begin{Highlighting}[]
\NormalTok{significant\_predictors }\OtherTok{\textless{}{-}}\NormalTok{ broom}\SpecialCharTok{::}\FunctionTok{tidy}\NormalTok{(logit\_model) }\SpecialCharTok{|\textgreater{}}
  \FunctionTok{arrange}\NormalTok{(p.value)}
\NormalTok{significant\_predictors}
\end{Highlighting}
\end{Shaded}

\begin{verbatim}
## # A tibble: 35 x 5
##    term                    estimate std.error statistic  p.value
##    <chr>                      <dbl>     <dbl>     <dbl>    <dbl>
##  1 work_interfereSometimes   2.69      0.296       9.07 1.16e-19
##  2 work_interfereOften       2.89      0.382       7.56 3.94e-14
##  3 work_interfereRarely      2.11      0.324       6.53 6.75e-11
##  4 family_historyYes         1.08      0.213       5.09 3.56e- 7
##  5 (Intercept)              -3.06      0.898      -3.40 6.63e- 4
##  6 GenderMale               -0.797     0.280      -2.84 4.45e- 3
##  7 coworkersYes              1.09      0.423       2.58 9.87e- 3
##  8 mental_careTRUE           0.578     0.229       2.52 1.17e- 2
##  9 Age                       0.0301    0.0145      2.08 3.78e- 2
## 10 no_employees26-100        0.658     0.364       1.81 7.04e- 2
## # i 25 more rows
\end{verbatim}

\begin{center}\rule{0.5\linewidth}{0.5pt}\end{center}

\paragraph{Feature Importance Plot}\label{feature-importance-plot}

\begin{Shaded}
\begin{Highlighting}[]
\NormalTok{feature\_importance }\OtherTok{\textless{}{-}}\NormalTok{ broom}\SpecialCharTok{::}\FunctionTok{tidy}\NormalTok{(logit\_model) }\SpecialCharTok{\%\textgreater{}\%}
  \FunctionTok{filter}\NormalTok{(term }\SpecialCharTok{!=} \StringTok{"(Intercept)"}\NormalTok{) }\SpecialCharTok{\%\textgreater{}\%}
  \FunctionTok{mutate}\NormalTok{(}\AttributeTok{odds\_ratio=}\FunctionTok{exp}\NormalTok{(estimate),}
         \AttributeTok{direction=}\FunctionTok{ifelse}\NormalTok{(estimate}\SpecialCharTok{\textgreater{}}\DecValTok{0}\NormalTok{,}\StringTok{"Increased"}\NormalTok{,}\StringTok{"Decreased"}\NormalTok{)) }\SpecialCharTok{\%\textgreater{}\%}
  \FunctionTok{arrange}\NormalTok{(}\FunctionTok{desc}\NormalTok{(}\FunctionTok{abs}\NormalTok{(estimate)))}
\NormalTok{top\_features }\OtherTok{\textless{}{-}}\NormalTok{ feature\_importance }\SpecialCharTok{\%\textgreater{}\%} \FunctionTok{head}\NormalTok{(}\DecValTok{15}\NormalTok{)}

\FunctionTok{ggplot}\NormalTok{(top\_features, }\FunctionTok{aes}\NormalTok{(}\AttributeTok{x=}\FunctionTok{reorder}\NormalTok{(term, estimate), }\AttributeTok{y=}\NormalTok{estimate, }\AttributeTok{fill=}\NormalTok{direction)) }\SpecialCharTok{+}
  \FunctionTok{geom\_col}\NormalTok{() }\SpecialCharTok{+}
  \FunctionTok{coord\_flip}\NormalTok{() }\SpecialCharTok{+}
  \FunctionTok{labs}\NormalTok{(}\AttributeTok{title=}\StringTok{"Top 15 Most Important Predictors of Mental Health Treatment Need"}\NormalTok{,}
       \AttributeTok{x=}\StringTok{"Predictor"}\NormalTok{, }\AttributeTok{y=}\StringTok{"Estimate (log{-}odds)"}\NormalTok{, }\AttributeTok{fill=}\StringTok{"Direction"}\NormalTok{) }\SpecialCharTok{+}
  \FunctionTok{theme\_minimal}\NormalTok{()}
\end{Highlighting}
\end{Shaded}

\pandocbounded{\includegraphics[keepaspectratio]{2501report_ShenHongshan_Version2_files/figure-latex/unnamed-chunk-13-1.pdf}}

\subsubsection{4.4 Insights}\label{insights-1}

\begin{itemize}
\item
  \textbf{Work interference} is the single strongest signal: respondents
  saying their mental health often interferes with work are far more
  likely to be in treatment. This supports screening based on
  self‑reported work interference.
\item
  \textbf{Family history} and \textbf{coworker openness} are also
  meaningful predictors. Family history increases risk, and feeling free
  to discuss with coworkers is positively associated with being in
  treatment, reflecting greater willingness to seek help or workplaces
  where issues are more openly discussed.
\item
  Company provisions (\textbf{mental\_care TRUE, anonymity TRUE}) are
  positively associated with treatment.This can mean:

  \begin{itemize}
  \item
    access to care increases uptake
  \item
    companies experiencing higher demand put programs in place (reverse
    causality)
  \item
    confounding by unmeasured factors.
  \end{itemize}
\item
  Male gender is associated with lower odds of being in treatment. That
  likely reflects \textbf{differing help‑seeking behavior}, \textbf{not
  necessarily lower need.}
\end{itemize}

\subsubsection{5. Model Evaluation}\label{model-evaluation}

\paragraph{ROC Curve}\label{roc-curve}

\begin{itemize}
\item
  ROC plots \textbf{sensitivity (true positive rate)} vs \textbf{1 −
  specificity (false positive rate)} over thresholds. The curve is
  \textbf{well above the diagonal:} model discriminates much better than
  random.
\item
  \textbf{AUC = 0.8476716} indicates good ranking/discrimination: the
  model usually gives higher predicted probabilities to respondents who
  are actually in treatment than to those who are not.
\end{itemize}

\begin{Shaded}
\begin{Highlighting}[]
\NormalTok{test\_data}\SpecialCharTok{$}\NormalTok{predicted\_prob }\OtherTok{\textless{}{-}} \FunctionTok{predict}\NormalTok{(logit\_model, }\AttributeTok{newdata=}\NormalTok{test\_data, }\AttributeTok{type=}\StringTok{"response"}\NormalTok{)}
\NormalTok{test\_data}\SpecialCharTok{$}\NormalTok{predicted\_class }\OtherTok{\textless{}{-}} \FunctionTok{ifelse}\NormalTok{(test\_data}\SpecialCharTok{$}\NormalTok{predicted\_prob }\SpecialCharTok{\textgreater{}} \FloatTok{0.5}\NormalTok{, }\DecValTok{1}\NormalTok{, }\DecValTok{0}\NormalTok{)}

\NormalTok{test\_conf }\OtherTok{\textless{}{-}} \FunctionTok{confusionMatrix}\NormalTok{(}\FunctionTok{as.factor}\NormalTok{(test\_data}\SpecialCharTok{$}\NormalTok{predicted\_class), }\FunctionTok{as.factor}\NormalTok{(test\_data}\SpecialCharTok{$}\NormalTok{treatment\_binary))}
\FunctionTok{print}\NormalTok{(test\_conf)}
\end{Highlighting}
\end{Shaded}

\begin{verbatim}
## Confusion Matrix and Statistics
## 
##           Reference
## Prediction   0   1
##          0  69  15
##          1  44 168
##                                           
##                Accuracy : 0.8007          
##                  95% CI : (0.7506, 0.8447)
##     No Information Rate : 0.6182          
##     P-Value [Acc > NIR] : 1.045e-11       
##                                           
##                   Kappa : 0.5559          
##                                           
##  Mcnemar's Test P-Value : 0.0002671       
##                                           
##             Sensitivity : 0.6106          
##             Specificity : 0.9180          
##          Pos Pred Value : 0.8214          
##          Neg Pred Value : 0.7925          
##              Prevalence : 0.3818          
##          Detection Rate : 0.2331          
##    Detection Prevalence : 0.2838          
##       Balanced Accuracy : 0.7643          
##                                           
##        'Positive' Class : 0               
## 
\end{verbatim}

\begin{Shaded}
\begin{Highlighting}[]
\NormalTok{roc\_curve }\OtherTok{\textless{}{-}} \FunctionTok{roc}\NormalTok{(test\_data}\SpecialCharTok{$}\NormalTok{treatment\_binary, test\_data}\SpecialCharTok{$}\NormalTok{predicted\_prob)}
\FunctionTok{cat}\NormalTok{(}\StringTok{"AUC on test:"}\NormalTok{, }\FunctionTok{auc}\NormalTok{(roc\_curve), }\StringTok{"}\SpecialCharTok{\textbackslash{}n}\StringTok{"}\NormalTok{)}
\end{Highlighting}
\end{Shaded}

\begin{verbatim}
## AUC on test: 0.8476716
\end{verbatim}

\begin{Shaded}
\begin{Highlighting}[]
\FunctionTok{plot}\NormalTok{(roc\_curve, }\AttributeTok{main=}\StringTok{"ROC Curve for Mental Health Treatment Prediction"}\NormalTok{, }\AttributeTok{col=}\StringTok{"blue"}\NormalTok{, }\AttributeTok{lwd=}\DecValTok{3}\NormalTok{)}
\FunctionTok{abline}\NormalTok{(}\AttributeTok{a=}\DecValTok{0}\NormalTok{, }\AttributeTok{b=}\DecValTok{1}\NormalTok{, }\AttributeTok{col=}\StringTok{"grey"}\NormalTok{, }\AttributeTok{lty=}\DecValTok{2}\NormalTok{)}
\end{Highlighting}
\end{Shaded}

\pandocbounded{\includegraphics[keepaspectratio]{2501report_ShenHongshan_Version2_files/figure-latex/unnamed-chunk-14-1.pdf}}

\paragraph{Prediction (training set)}\label{prediction-training-set}

\begin{itemize}
\item
  Strength of model:\\
  \strut ~\textbf{High Specificity (91.91\%)}: Excellent at identifying
  people who truly need treatment\\
  \strut ~\textbf{Good Overall Accuracy (80.23\%)}: Beats baseline by 16
  percentage points\\
  \strut ~\textbf{Balanced Precision}: Similar precision for both
  classes (\textasciitilde80\%)\\
  \strut ~\textbf{Statistical Significance}: Results are not due to
  chance
\item
  Concerns\\
  \strut ~\textbf{Low Sensitivity (59.27\%)}: Poor at identifying people
  who don't need treatment\\
  \strut ~\textbf{Imbalanced Performance}: Much better at detecting
  treatment needs than non-needs\\
  \strut ~\textbf{Potential Over-intervention}: 40.73\% of ``No
  treatment'' people are incorrectly flagged
\end{itemize}

\begin{Shaded}
\begin{Highlighting}[]
\NormalTok{train\_data}\SpecialCharTok{$}\NormalTok{predicted\_prob }\OtherTok{\textless{}{-}} \FunctionTok{predict}\NormalTok{(logit\_model, }\AttributeTok{type=}\StringTok{"response"}\NormalTok{)}
\NormalTok{train\_data}\SpecialCharTok{$}\NormalTok{predicted\_class }\OtherTok{\textless{}{-}} \FunctionTok{ifelse}\NormalTok{(train\_data}\SpecialCharTok{$}\NormalTok{predicted\_prob }\SpecialCharTok{\textgreater{}} \FloatTok{0.5}\NormalTok{, }\DecValTok{1}\NormalTok{, }\DecValTok{0}\NormalTok{)}
\NormalTok{train\_conf }\OtherTok{\textless{}{-}} \FunctionTok{confusionMatrix}\NormalTok{(}\FunctionTok{as.factor}\NormalTok{(train\_data}\SpecialCharTok{$}\NormalTok{predicted\_class), }\FunctionTok{as.factor}\NormalTok{(train\_data}\SpecialCharTok{$}\NormalTok{treatment\_binary))}
\FunctionTok{print}\NormalTok{(train\_conf)}
\end{Highlighting}
\end{Shaded}

\begin{verbatim}
## Confusion Matrix and Statistics
## 
##           Reference
## Prediction   0   1
##          0 147  36
##          1 101 409
##                                           
##                Accuracy : 0.8023          
##                  95% CI : (0.7707, 0.8313)
##     No Information Rate : 0.6421          
##     P-Value [Acc > NIR] : < 2.2e-16       
##                                           
##                   Kappa : 0.5434          
##                                           
##  Mcnemar's Test P-Value : 4.554e-08       
##                                           
##             Sensitivity : 0.5927          
##             Specificity : 0.9191          
##          Pos Pred Value : 0.8033          
##          Neg Pred Value : 0.8020          
##              Prevalence : 0.3579          
##          Detection Rate : 0.2121          
##    Detection Prevalence : 0.2641          
##       Balanced Accuracy : 0.7559          
##                                           
##        'Positive' Class : 0               
## 
\end{verbatim}

\begin{Shaded}
\begin{Highlighting}[]
\CommentTok{\# vusualize confusion data}
\NormalTok{conf\_matrix\_data }\OtherTok{\textless{}{-}} \FunctionTok{as.data.frame}\NormalTok{(}\FunctionTok{as.table}\NormalTok{(train\_conf}\SpecialCharTok{$}\NormalTok{table))}
\FunctionTok{names}\NormalTok{(conf\_matrix\_data) }\OtherTok{\textless{}{-}} \FunctionTok{c}\NormalTok{(}\StringTok{"Predicted"}\NormalTok{, }\StringTok{"Actual"}\NormalTok{, }\StringTok{"Freq"}\NormalTok{)}

\CommentTok{\# calculate relevant statistics}
\NormalTok{conf\_matrix\_data }\OtherTok{\textless{}{-}}\NormalTok{ conf\_matrix\_data }\SpecialCharTok{\%\textgreater{}\%}
  \FunctionTok{group\_by}\NormalTok{(Actual) }\SpecialCharTok{\%\textgreater{}\%}
  \FunctionTok{mutate}\NormalTok{(}
    \AttributeTok{Percentage =}\NormalTok{ Freq }\SpecialCharTok{/} \FunctionTok{sum}\NormalTok{(Freq) }\SpecialCharTok{*} \DecValTok{100}\NormalTok{,}
    \AttributeTok{Label =} \FunctionTok{paste0}\NormalTok{(Freq, }\StringTok{"}\SpecialCharTok{\textbackslash{}n}\StringTok{("}\NormalTok{, }\FunctionTok{round}\NormalTok{(Percentage, }\DecValTok{1}\NormalTok{), }\StringTok{"\%)"}\NormalTok{)}
\NormalTok{  )}

\CommentTok{\# plot heatmap ofr confusion matrix}
\FunctionTok{ggplot}\NormalTok{(conf\_matrix\_data, }\FunctionTok{aes}\NormalTok{(}\AttributeTok{x =}\NormalTok{ Actual, }\AttributeTok{y =}\NormalTok{ Predicted, }\AttributeTok{fill =}\NormalTok{ Freq)) }\SpecialCharTok{+}
  \FunctionTok{geom\_tile}\NormalTok{(}\AttributeTok{color =} \StringTok{"white"}\NormalTok{, }\AttributeTok{alpha =} \FloatTok{0.8}\NormalTok{) }\SpecialCharTok{+}
  \FunctionTok{geom\_text}\NormalTok{(}\FunctionTok{aes}\NormalTok{(}\AttributeTok{label =}\NormalTok{ Label), }\AttributeTok{color =} \StringTok{"black"}\NormalTok{, }\AttributeTok{size =} \DecValTok{5}\NormalTok{) }\SpecialCharTok{+}
  \FunctionTok{scale\_fill\_gradient}\NormalTok{(}\AttributeTok{low =} \StringTok{"\#F8F9FA"}\NormalTok{, }\AttributeTok{high =} \StringTok{"blue"}\NormalTok{, }\AttributeTok{name =} \StringTok{"frequency"}\NormalTok{) }\SpecialCharTok{+}
  \FunctionTok{labs}\NormalTok{(}
    \AttributeTok{title =} \StringTok{"heatmap ofr confusion matrix"}\NormalTok{,}
    \AttributeTok{subtitle =} \FunctionTok{paste}\NormalTok{(}\StringTok{"Accuracy:"}\NormalTok{, }\FunctionTok{round}\NormalTok{(train\_conf}\SpecialCharTok{$}\NormalTok{overall[}\StringTok{"Accuracy"}\NormalTok{]}\SpecialCharTok{*}\DecValTok{100}\NormalTok{, }\DecValTok{1}\NormalTok{), }\StringTok{"\%"}\NormalTok{),}
    \AttributeTok{x =} \StringTok{"Actural class"}\NormalTok{,}
    \AttributeTok{y =} \StringTok{"Predicted class"}\NormalTok{,}
    \AttributeTok{caption =} \FunctionTok{paste}\NormalTok{(}\StringTok{"sensitivity:"}\NormalTok{, }\FunctionTok{round}\NormalTok{(train\_conf}\SpecialCharTok{$}\NormalTok{byClass[}\StringTok{"Sensitivity"}\NormalTok{]}\SpecialCharTok{*}\DecValTok{100}\NormalTok{, }\DecValTok{1}\NormalTok{), }\StringTok{"\% |"}\NormalTok{,}
                   \StringTok{"Specificity:"}\NormalTok{, }\FunctionTok{round}\NormalTok{(train\_conf}\SpecialCharTok{$}\NormalTok{byClass[}\StringTok{"Specificity"}\NormalTok{]}\SpecialCharTok{*}\DecValTok{100}\NormalTok{, }\DecValTok{1}\NormalTok{), }\StringTok{"\%"}\NormalTok{)}
\NormalTok{  ) }\SpecialCharTok{+}
  \FunctionTok{theme\_minimal}\NormalTok{() }
\end{Highlighting}
\end{Shaded}

\pandocbounded{\includegraphics[keepaspectratio]{2501report_ShenHongshan_Version2_files/figure-latex/unnamed-chunk-15-1.pdf}}

\subsection{\texorpdfstring{\textbf{6. Answering data science
questions}}{6. Answering data science questions}}\label{answering-data-science-questions}

\subsubsection{6.1. How to identify potential mental issue, In other
words, what features are most predictive of needing mental health
treatment?}\label{how-to-identify-potential-mental-issue-in-other-words-what-features-are-most-predictive-of-needing-mental-health-treatment}

\begin{itemize}
\tightlist
\item
  We identified that \textbf{family history}, \textbf{work
  interference}, and \textbf{workplace support systems} are among the
  strongest predictors of mental health treatment need. (detailed
  analysis seen in \textbf{section 3.4 \& 4.4})
\end{itemize}

\subsubsection{6.2 What could employers do to reduce mental issue's
risks?}\label{what-could-employers-do-to-reduce-mental-issues-risks}

\begin{itemize}
\tightlist
\item
  A logistic regression model giving \textbf{the most predictive
  factors} can reasonably \textbf{identify potnetial employees at higher
  risk.}
\item
  Companies should focus on improving \textbf{support policies} for
  employees showing risk factors.

  \begin{itemize}
  \tightlist
  \item
    easier access to a\textbf{pplying for medical leave due to mental
    issue}
  \item
    protect workers \textbf{anonymity regarding mental health issue},
    and creating a non-judgmental environment
  \end{itemize}
\item
  Increasing mental health care benefits, easy-access mental care
  treatment resources.

  \begin{itemize}
  \tightlist
  \item
    \textbf{interesting} to notice the \textbf{different approach} that
    may apply to different Gender (see visualization below!)

    \begin{itemize}
    \tightlist
    \item
      \textbf{women are more sensitive to mental care benefit (money)}:
      providing more money on mental health care have a greater positive
      effect on seeking treatment for women(nearly \textbf{30\%} more)
      compared to men(around \textbf{10\%} more).
    \item
      on the other hand, \textbf{men are more sensitive to giving health
      care access}: providing more money on mental health care have a
      greater positive effect on seeking treatment for men(around
      \textbf{15\%} increase) compared to women(lower than \textbf{5\%}
      increase).
    \item
      giving \textbf{direct mental care options at work place}, is the
      \textbf{most effective} way to encourage employees seeking
      treatment, for both men and women.
    \end{itemize}
  \end{itemize}
\end{itemize}

\begin{Shaded}
\begin{Highlighting}[]
\CommentTok{\# 1. Interaction Effects Visualization {-} benefit * Gender}
\NormalTok{gender\_care\_data }\OtherTok{\textless{}{-}}\NormalTok{ survey\_clean0 }\SpecialCharTok{|\textgreater{}}
  \FunctionTok{filter}\NormalTok{(Gender}\SpecialCharTok{!=}\StringTok{"Transgender/Non{-}binary/Others"}\NormalTok{)}\SpecialCharTok{|\textgreater{}}
  \FunctionTok{mutate}\NormalTok{(}
    \AttributeTok{treatment\_binary =} \FunctionTok{ifelse}\NormalTok{(treatment }\SpecialCharTok{==} \StringTok{"Yes"}\NormalTok{, }\DecValTok{1}\NormalTok{, }\DecValTok{0}\NormalTok{)}
\NormalTok{  )}\SpecialCharTok{|\textgreater{}}
  \FunctionTok{group\_by}\NormalTok{(Gender, benefits) }\SpecialCharTok{|\textgreater{}}
  \FunctionTok{summarise}\NormalTok{(}
    \AttributeTok{treatment\_rate =} \FunctionTok{mean}\NormalTok{(treatment\_binary),}
    \AttributeTok{n =} \FunctionTok{n}\NormalTok{(),}
    \AttributeTok{.groups =} \StringTok{\textquotesingle{}drop\textquotesingle{}}
\NormalTok{  )}

\FunctionTok{ggplot}\NormalTok{(gender\_care\_data, }\FunctionTok{aes}\NormalTok{(}\AttributeTok{x =}\NormalTok{ Gender, }\AttributeTok{y =}\NormalTok{ treatment\_rate, }\AttributeTok{fill =}\NormalTok{ benefits)) }\SpecialCharTok{+}
  \FunctionTok{geom\_bar}\NormalTok{(}\AttributeTok{stat =} \StringTok{"identity"}\NormalTok{, }\AttributeTok{position =} \StringTok{"dodge"}\NormalTok{) }\SpecialCharTok{+}
  \FunctionTok{geom\_text}\NormalTok{(}\FunctionTok{aes}\NormalTok{(}\AttributeTok{label =} \FunctionTok{paste0}\NormalTok{(}\FunctionTok{round}\NormalTok{(treatment\_rate}\SpecialCharTok{*}\DecValTok{100}\NormalTok{, }\DecValTok{1}\NormalTok{), }\StringTok{"\%"}\NormalTok{)), }
            \AttributeTok{position =} \FunctionTok{position\_dodge}\NormalTok{(}\AttributeTok{width =} \FloatTok{0.9}\NormalTok{), }\AttributeTok{vjust =} \SpecialCharTok{{-}}\FloatTok{0.5}\NormalTok{) }\SpecialCharTok{+}
  \FunctionTok{labs}\NormalTok{(}
    \AttributeTok{title =} \StringTok{"Interaction Effect: Gender × mental health benefits"}\NormalTok{,}
    \AttributeTok{x =} \StringTok{"Gender"}\NormalTok{,}
    \AttributeTok{y =} \StringTok{"Treatment Rate"}\NormalTok{,}
    \AttributeTok{fill =} \StringTok{"Company Provides}\SpecialCharTok{\textbackslash{}n}\StringTok{Mental Health Benefits"}
\NormalTok{  ) }\SpecialCharTok{+}
  \FunctionTok{theme\_minimal}\NormalTok{() }\SpecialCharTok{+}
  \FunctionTok{scale\_y\_continuous}\NormalTok{(}\AttributeTok{labels =}\NormalTok{ scales}\SpecialCharTok{::}\NormalTok{percent) }\SpecialCharTok{+}
  \FunctionTok{scale\_fill\_brewer}\NormalTok{(}\AttributeTok{palette =} \StringTok{"Set2"}\NormalTok{)}
\end{Highlighting}
\end{Shaded}

\pandocbounded{\includegraphics[keepaspectratio]{2501report_ShenHongshan_Version2_files/figure-latex/unnamed-chunk-16-1.pdf}}

\begin{Shaded}
\begin{Highlighting}[]
\CommentTok{\# 2. Interaction Effects Visualization {-} seek\_help * Gender}
\NormalTok{gender\_care\_data }\OtherTok{\textless{}{-}}\NormalTok{ survey\_clean0 }\SpecialCharTok{|\textgreater{}}
  \FunctionTok{filter}\NormalTok{(Gender}\SpecialCharTok{!=}\StringTok{"Transgender/Non{-}binary/Others"}\NormalTok{)}\SpecialCharTok{|\textgreater{}}
  \FunctionTok{mutate}\NormalTok{(}
    \AttributeTok{treatment\_binary =} \FunctionTok{ifelse}\NormalTok{(treatment }\SpecialCharTok{==} \StringTok{"Yes"}\NormalTok{, }\DecValTok{1}\NormalTok{, }\DecValTok{0}\NormalTok{)}
\NormalTok{  )}\SpecialCharTok{|\textgreater{}}
  \FunctionTok{group\_by}\NormalTok{(Gender, seek\_help) }\SpecialCharTok{|\textgreater{}}
  \FunctionTok{summarise}\NormalTok{(}
    \AttributeTok{treatment\_rate =} \FunctionTok{mean}\NormalTok{(treatment\_binary),}
    \AttributeTok{n =} \FunctionTok{n}\NormalTok{(),}
    \AttributeTok{.groups =} \StringTok{\textquotesingle{}drop\textquotesingle{}}
\NormalTok{  )}

\FunctionTok{ggplot}\NormalTok{(gender\_care\_data, }\FunctionTok{aes}\NormalTok{(}\AttributeTok{x =}\NormalTok{ Gender, }\AttributeTok{y =}\NormalTok{ treatment\_rate, }\AttributeTok{fill =}\NormalTok{ seek\_help)) }\SpecialCharTok{+}
  \FunctionTok{geom\_bar}\NormalTok{(}\AttributeTok{stat =} \StringTok{"identity"}\NormalTok{, }\AttributeTok{position =} \StringTok{"dodge"}\NormalTok{) }\SpecialCharTok{+}
  \FunctionTok{geom\_text}\NormalTok{(}\FunctionTok{aes}\NormalTok{(}\AttributeTok{label =} \FunctionTok{paste0}\NormalTok{(}\FunctionTok{round}\NormalTok{(treatment\_rate}\SpecialCharTok{*}\DecValTok{100}\NormalTok{, }\DecValTok{1}\NormalTok{), }\StringTok{"\%"}\NormalTok{)), }
            \AttributeTok{position =} \FunctionTok{position\_dodge}\NormalTok{(}\AttributeTok{width =} \FloatTok{0.9}\NormalTok{), }\AttributeTok{vjust =} \SpecialCharTok{{-}}\FloatTok{0.5}\NormalTok{) }\SpecialCharTok{+}
  \FunctionTok{labs}\NormalTok{(}
    \AttributeTok{title =} \StringTok{"Interaction Effect: Gender × Access to help"}\NormalTok{,}
    \AttributeTok{x =} \StringTok{"Gender"}\NormalTok{,}
    \AttributeTok{y =} \StringTok{"Treatment Rate"}\NormalTok{,}
    \AttributeTok{fill =} \StringTok{"Company Provides}\SpecialCharTok{\textbackslash{}n}\StringTok{Access to help"}
\NormalTok{  ) }\SpecialCharTok{+}
  \FunctionTok{theme\_minimal}\NormalTok{() }\SpecialCharTok{+}
  \FunctionTok{scale\_y\_continuous}\NormalTok{(}\AttributeTok{labels =}\NormalTok{ scales}\SpecialCharTok{::}\NormalTok{percent) }\SpecialCharTok{+}
  \FunctionTok{scale\_fill\_brewer}\NormalTok{(}\AttributeTok{palette =} \StringTok{"Set2"}\NormalTok{)}
\end{Highlighting}
\end{Shaded}

\pandocbounded{\includegraphics[keepaspectratio]{2501report_ShenHongshan_Version2_files/figure-latex/unnamed-chunk-16-2.pdf}}

\begin{Shaded}
\begin{Highlighting}[]
\CommentTok{\# 3. Interaction Effects Visualization {-} care\_options * Gender}
\NormalTok{gender\_care\_data }\OtherTok{\textless{}{-}}\NormalTok{ survey\_clean0 }\SpecialCharTok{|\textgreater{}}
  \FunctionTok{filter}\NormalTok{(Gender}\SpecialCharTok{!=}\StringTok{"Transgender/Non{-}binary/Others"}\NormalTok{)}\SpecialCharTok{|\textgreater{}}
  \FunctionTok{mutate}\NormalTok{(}
    \AttributeTok{treatment\_binary =} \FunctionTok{ifelse}\NormalTok{(treatment }\SpecialCharTok{==} \StringTok{"Yes"}\NormalTok{, }\DecValTok{1}\NormalTok{, }\DecValTok{0}\NormalTok{)}
\NormalTok{  )}\SpecialCharTok{|\textgreater{}}
  \FunctionTok{group\_by}\NormalTok{(Gender, care\_options) }\SpecialCharTok{|\textgreater{}}
  \FunctionTok{summarise}\NormalTok{(}
    \AttributeTok{treatment\_rate =} \FunctionTok{mean}\NormalTok{(treatment\_binary),}
    \AttributeTok{n =} \FunctionTok{n}\NormalTok{(),}
    \AttributeTok{.groups =} \StringTok{\textquotesingle{}drop\textquotesingle{}}
\NormalTok{  )}

\FunctionTok{ggplot}\NormalTok{(gender\_care\_data, }\FunctionTok{aes}\NormalTok{(}\AttributeTok{x =}\NormalTok{ Gender, }\AttributeTok{y =}\NormalTok{ treatment\_rate, }\AttributeTok{fill =}\NormalTok{ care\_options)) }\SpecialCharTok{+}
  \FunctionTok{geom\_bar}\NormalTok{(}\AttributeTok{stat =} \StringTok{"identity"}\NormalTok{, }\AttributeTok{position =} \StringTok{"dodge"}\NormalTok{) }\SpecialCharTok{+}
  \FunctionTok{geom\_text}\NormalTok{(}\FunctionTok{aes}\NormalTok{(}\AttributeTok{label =} \FunctionTok{paste0}\NormalTok{(}\FunctionTok{round}\NormalTok{(treatment\_rate}\SpecialCharTok{*}\DecValTok{100}\NormalTok{, }\DecValTok{1}\NormalTok{), }\StringTok{"\%"}\NormalTok{)), }
            \AttributeTok{position =} \FunctionTok{position\_dodge}\NormalTok{(}\AttributeTok{width =} \FloatTok{0.9}\NormalTok{), }\AttributeTok{vjust =} \SpecialCharTok{{-}}\FloatTok{0.5}\NormalTok{) }\SpecialCharTok{+}
  \FunctionTok{labs}\NormalTok{(}
    \AttributeTok{title =} \StringTok{"Interaction Effect: Gender × mental care options"}\NormalTok{,}
    \AttributeTok{x =} \StringTok{"Gender"}\NormalTok{,}
    \AttributeTok{y =} \StringTok{"Treatment Rate"}\NormalTok{,}
    \AttributeTok{fill =} \StringTok{"Company Provides}\SpecialCharTok{\textbackslash{}n}\StringTok{Mental care options"}
\NormalTok{  ) }\SpecialCharTok{+}
  \FunctionTok{theme\_minimal}\NormalTok{() }\SpecialCharTok{+}
  \FunctionTok{scale\_y\_continuous}\NormalTok{(}\AttributeTok{labels =}\NormalTok{ scales}\SpecialCharTok{::}\NormalTok{percent) }\SpecialCharTok{+}
  \FunctionTok{scale\_fill\_brewer}\NormalTok{(}\AttributeTok{palette =} \StringTok{"Set2"}\NormalTok{)}
\end{Highlighting}
\end{Shaded}

\pandocbounded{\includegraphics[keepaspectratio]{2501report_ShenHongshan_Version2_files/figure-latex/unnamed-chunk-16-3.pdf}}

\subsection{7. Limitation \& possible
improvement}\label{limitation-possible-improvement}

\begin{itemize}
\tightlist
\item
  \textbf{Most significant limitation:} Subjective survey responses,
  possible response bias.

  \begin{itemize}
  \tightlist
  \item
    \textbf{target variable ``treatment''} indicates whether respondents
    are currently receiving mental health treatment. It may demonstrate
    the actual need for mental treatment (\textbf{real mental issue
    posed}). However, due to subjective survey response
    \textbf{``willingness to get treatment''} also affected whether or
    not they are receiving treatment. This \textbf{self-report bias}
    means some individuals with mental health conditions may not seek or
    report treatment, while others receiving treatment may do so
    preventively rather than for diagnosed conditions.
  \end{itemize}
\item
  \textbf{one-to-one association} between one feature and target
  variable''treatment'' is investigated, but the \textbf{interactions
  between} different features and how their combinations is associated
  with need for treatment is not considered. this is an aspect of future
  investigation.

  \begin{itemize}
  \tightlist
  \item
    an example is given below to explore the relationship between two
    factors \textbf{``work\_interfere'' \& ``family\_history''}
  \item
    \textbf{Heatmap} indicating the correlation between features are
    provided below
  \end{itemize}
\item
  assessing the \textbf{association} between demographic\&workplace
  features and treatment is enough to \textbf{identify individuals at
  risk}, but from employers prospect, if they want to address mental
  issue for employees more effectively, analyzing the causual
  relationship should be considered.

  \begin{itemize}
  \tightlist
  \item
    This can be solved by using more advanced ML (random forest,
    boosting), longitudinal data, and testing interventions\ldots{}
  \end{itemize}
\end{itemize}

\subsubsection{example demonstration:}\label{example-demonstration}

\begin{enumerate}
\def\labelenumi{\arabic{enumi}.}
\tightlist
\item
  simple correlation to interaction
\end{enumerate}

\begin{Shaded}
\begin{Highlighting}[]
\CommentTok{\# 1. Interaction Effects Visualization {-} work\_interfere * family\_history}
\NormalTok{interaction\_data }\OtherTok{\textless{}{-}}\NormalTok{ train\_data }\SpecialCharTok{\%\textgreater{}\%}
  \FunctionTok{group\_by}\NormalTok{(work\_interfere, family\_history) }\SpecialCharTok{\%\textgreater{}\%}
  \FunctionTok{summarise}\NormalTok{(}
    \AttributeTok{treatment\_rate =} \FunctionTok{mean}\NormalTok{(treatment\_binary),}
    \AttributeTok{n =} \FunctionTok{n}\NormalTok{(),}
    \AttributeTok{.groups =} \StringTok{\textquotesingle{}drop\textquotesingle{}}
\NormalTok{  ) }\SpecialCharTok{\%\textgreater{}\%}
  \FunctionTok{filter}\NormalTok{(}\SpecialCharTok{!}\FunctionTok{is.na}\NormalTok{(work\_interfere))}

\FunctionTok{ggplot}\NormalTok{(interaction\_data, }\FunctionTok{aes}\NormalTok{(}\AttributeTok{x =}\NormalTok{ work\_interfere, }\AttributeTok{y =}\NormalTok{ treatment\_rate, }
                             \AttributeTok{color =}\NormalTok{ family\_history, }\AttributeTok{group =}\NormalTok{ family\_history)) }\SpecialCharTok{+}
  \FunctionTok{geom\_line}\NormalTok{(}\AttributeTok{size =} \FloatTok{1.5}\NormalTok{) }\SpecialCharTok{+}
  \FunctionTok{geom\_point}\NormalTok{(}\AttributeTok{size =} \DecValTok{3}\NormalTok{) }\SpecialCharTok{+}
  \FunctionTok{geom\_text}\NormalTok{(}\FunctionTok{aes}\NormalTok{(}\AttributeTok{label =} \FunctionTok{paste0}\NormalTok{(}\FunctionTok{round}\NormalTok{(treatment\_rate}\SpecialCharTok{*}\DecValTok{100}\NormalTok{, }\DecValTok{1}\NormalTok{), }\StringTok{"\%"}\NormalTok{)), }
            \AttributeTok{vjust =} \SpecialCharTok{{-}}\FloatTok{0.8}\NormalTok{, }\AttributeTok{size =} \DecValTok{3}\NormalTok{) }\SpecialCharTok{+}
  \FunctionTok{labs}\NormalTok{(}
    \AttributeTok{title =} \StringTok{"Interaction Effect: Work Interference × Family History"}\NormalTok{,}
    \AttributeTok{subtitle =} \StringTok{"Family history significantly amplifies the impact of work interference on treatment needs"}\NormalTok{,}
    \AttributeTok{x =} \StringTok{"Frequency of Work Interference"}\NormalTok{,}
    \AttributeTok{y =} \StringTok{"Treatment Rate"}\NormalTok{,}
    \AttributeTok{color =} \StringTok{"Family History of}\SpecialCharTok{\textbackslash{}n}\StringTok{Mental Health"}
\NormalTok{  ) }\SpecialCharTok{+}
  \FunctionTok{theme\_minimal}\NormalTok{() }\SpecialCharTok{+}
  \FunctionTok{scale\_y\_continuous}\NormalTok{(}\AttributeTok{labels =}\NormalTok{ scales}\SpecialCharTok{::}\NormalTok{percent) }\SpecialCharTok{+}
  \FunctionTok{scale\_color\_manual}\NormalTok{(}\AttributeTok{values =} \FunctionTok{c}\NormalTok{(}\StringTok{"No"} \OtherTok{=} \StringTok{"red"}\NormalTok{, }\StringTok{"Yes"} \OtherTok{=} \StringTok{"blue"}\NormalTok{))}
\end{Highlighting}
\end{Shaded}

\pandocbounded{\includegraphics[keepaspectratio]{2501report_ShenHongshan_Version2_files/figure-latex/unnamed-chunk-17-1.pdf}}

\begin{itemize}
\item
  There is a powerful interaction between work interference and family
  history:

  \begin{itemize}
  \item
    For employees without family history, even frequent work
    interference leads to only 60\% treatment rate.
  \item
    But for those with family history, the same level of work
    interference drives treatment rates above 90\%.
  \end{itemize}
\item
  This tells us screening should \textbf{prioritize employees who report
  both family history AND work interference.}
\end{itemize}

\begin{enumerate}
\def\labelenumi{\arabic{enumi}.}
\setcounter{enumi}{1}
\tightlist
\item
  correlation heatmap
\end{enumerate}

\begin{Shaded}
\begin{Highlighting}[]
\ControlFlowTok{if}\NormalTok{(}\SpecialCharTok{!}\FunctionTok{requireNamespace}\NormalTok{(}\StringTok{"corrplot"}\NormalTok{, }\AttributeTok{quietly=}\ConstantTok{TRUE}\NormalTok{)) }\FunctionTok{install.packages}\NormalTok{(}\StringTok{"corrplot"}\NormalTok{)}
\ControlFlowTok{if}\NormalTok{(}\SpecialCharTok{!}\FunctionTok{requireNamespace}\NormalTok{(}\StringTok{"reshape2"}\NormalTok{, }\AttributeTok{quietly=}\ConstantTok{TRUE}\NormalTok{)) }\FunctionTok{install.packages}\NormalTok{(}\StringTok{"reshape2"}\NormalTok{)}
\ControlFlowTok{if}\NormalTok{(}\SpecialCharTok{!}\FunctionTok{requireNamespace}\NormalTok{(}\StringTok{"Hmisc"}\NormalTok{, }\AttributeTok{quietly=}\ConstantTok{TRUE}\NormalTok{)) }\FunctionTok{install.packages}\NormalTok{(}\StringTok{"Hmisc"}\NormalTok{)}
\FunctionTok{library}\NormalTok{(corrplot); }\FunctionTok{library}\NormalTok{(reshape2); }\FunctionTok{library}\NormalTok{(Hmisc)}

\NormalTok{corr\_df }\OtherTok{\textless{}{-}}\NormalTok{ model\_data }

\NormalTok{corr\_num }\OtherTok{\textless{}{-}}\NormalTok{ corr\_df }\SpecialCharTok{|\textgreater{}}
  \FunctionTok{mutate}\NormalTok{(}\FunctionTok{across}\NormalTok{(}\FunctionTok{everything}\NormalTok{(), }\SpecialCharTok{\textasciitilde{}} \ControlFlowTok{if}\NormalTok{(}\FunctionTok{is.factor}\NormalTok{(.)) }\FunctionTok{as.numeric}\NormalTok{(}\FunctionTok{as.factor}\NormalTok{(.)) }\ControlFlowTok{else}\NormalTok{ .))}

\NormalTok{corr\_mat }\OtherTok{\textless{}{-}} \FunctionTok{cor}\NormalTok{(corr\_num, }\AttributeTok{use =} \StringTok{"pairwise.complete.obs"}\NormalTok{, }\AttributeTok{method =} \StringTok{"spearman"}\NormalTok{)}

\NormalTok{corr\_mat[}\FunctionTok{upper.tri}\NormalTok{(corr\_mat, }\AttributeTok{diag =} \ConstantTok{TRUE}\NormalTok{)] }\OtherTok{\textless{}{-}} \ConstantTok{NA}

\NormalTok{melted }\OtherTok{\textless{}{-}}\NormalTok{ reshape2}\SpecialCharTok{::}\FunctionTok{melt}\NormalTok{(corr\_mat, }\AttributeTok{na.rm =} \ConstantTok{TRUE}\NormalTok{)}
\FunctionTok{ggplot}\NormalTok{(melted, }\FunctionTok{aes}\NormalTok{(}\AttributeTok{x =}\NormalTok{ Var2, }\AttributeTok{y =}\NormalTok{ Var1, }\AttributeTok{fill =}\NormalTok{ value)) }\SpecialCharTok{+}
  \FunctionTok{geom\_tile}\NormalTok{(}\AttributeTok{color =} \StringTok{"white"}\NormalTok{) }\SpecialCharTok{+}
  \FunctionTok{scale\_fill\_gradient2}\NormalTok{(}\AttributeTok{low =} \StringTok{"blue"}\NormalTok{, }\AttributeTok{mid =} \StringTok{"white"}\NormalTok{, }\AttributeTok{high =} \StringTok{"red"}\NormalTok{, }\AttributeTok{midpoint =} \DecValTok{0}\NormalTok{, }\AttributeTok{limits =} \FunctionTok{c}\NormalTok{(}\SpecialCharTok{{-}}\FloatTok{0.5}\NormalTok{,}\FloatTok{0.6}\NormalTok{)) }\SpecialCharTok{+}
  \FunctionTok{theme\_minimal}\NormalTok{() }\SpecialCharTok{+}
  \FunctionTok{theme}\NormalTok{(}\AttributeTok{axis.text.x =} \FunctionTok{element\_text}\NormalTok{(}\AttributeTok{angle =} \DecValTok{45}\NormalTok{, }\AttributeTok{hjust =} \DecValTok{1}\NormalTok{)) }\SpecialCharTok{+}
  \FunctionTok{labs}\NormalTok{(}\AttributeTok{x =} \StringTok{""}\NormalTok{, }\AttributeTok{y =} \StringTok{""}\NormalTok{, }\AttributeTok{fill =} \StringTok{"Spearman}\SpecialCharTok{\textbackslash{}n}\StringTok{rho"}\NormalTok{,}
       \AttributeTok{title =} \StringTok{"Spearman Correlation Heatmap (factors {-}\textgreater{} numeric)"}\NormalTok{)}
\end{Highlighting}
\end{Shaded}

\pandocbounded{\includegraphics[keepaspectratio]{2501report_ShenHongshan_Version2_files/figure-latex/correlation-heatmap-1.pdf}}

\begin{itemize}
\item
  notice the correlation between \textbf{supervisor} (feel free to
  discuss mental issue with direct supervisor) \& \textbf{coworkers
  (}feel free to discuss mental issue with coworkers) is quite high

  \begin{itemize}
  \item
    is there any potential causation ? any insightful management
    approach for employers?
  \item
    Supervisors may serve as \textbf{tone-setters} for team culture.
    When supervisors openly discuss mental health and demonstrate
    empathy, they model vulnerable communication, giving employees
    `permission' to have similar conversations with peers. This aligns
    with social learning theory, where employees observe and mimic
    leader behaviors.
  \item
    A supportive peer culture might \textbf{normalize} mental health
    discussions, making it safer for employees to approach supervisors.
    This creates a \textbf{virtuous cycle} where peer support builds
    confidence for upward communication.
  \end{itemize}
\end{itemize}

\subsection{8. Conclusion}\label{conclusion}

This study utilized the OSMI Mental Health in Tech Survey to investigate
the relationship between workplace factors, demographic characteristics,
and mental health treatment-seeking behavior. My analysis revealed that
\textbf{work interference} and \textbf{family history} of mental illness
are the \textbf{strongest predictors} of treatment need, with employees
reporting frequent work interference and those with family history
showing substantially higher treatment rates. Additionally, company's
\textbf{support systems}---including mental health benefits,
confidential resources, and supportive workplace
cultures---significantly correlate with treatment engagement. The
logistic regression model demonstrated good predictive capability with
an AUC of 0.85 and overall accuracy of 80\%, though it exhibited higher
specificity than sensitivity, indicating better performance at
identifying treatment needs than non-needs.

From a practical perspective, these findings suggest employers should
implement \textbf{targeted screening for employees experiencing work
interference}, particularly those with \textbf{family mental health
history,} and develop \textbf{gender-sensitive interventions} given the
pronounced differences in treatment-seeking patterns. The strong
correlation between supervisor and coworker support highlights the
importance of creating psychologically safe environments at multiple
organizational levels. While these associative findings provide valuable
insights for identifying at-risk individuals, establishing causal
relationships would further enhance employers' ability to design
effective interventions. Future research employing longitudinal designs
and causal inference methods could build upon these findings to develop
more targeted workplace mental health strategies that promote employee
well-being and organizational resilience.

\subsection{reference:}\label{reference}

OSMI Mental Health in Tech Survey
\href{https://www.kaggle.com/datasets/osmi/mental-health-in-tech-survey/data}{Kaggle
Dataset} - \href{https://www.statlearning.com/}{Introduction to
Statistical Learning}

\end{document}
